\section{The analytic outer-carrier construction}\label{app:carrier}
This appendix supplies the analytic details of Theorem~\ref{thm:R0}.
Throughout put \(s=\St=q^*\). Proposition~\ref{prop:R1} gives
\(1/4<s<13/50<1/3\). The only external \(I_{3322}\) premises are
those in External premise~\ref{prem:Coladangelo}.

\subsection{Finite realization}\label{app:realization}
For a path in \eqref{eq:J}, on the standard basis of \(\CC^m\) define
diagonal operators by
\[
 ((Z_A)_{ii},(Z_B)_{ii})=
 \begin{cases}
  (\gamma_{i-1},\gamma_i),&i\text{ odd},\\
  (-\gamma_i,-\gamma_{i-1}),&i\text{ even}.
 \end{cases}
\]
On even edges \((i,i+1)\), let \(X_A\) be the swap with weight
\(w(\gamma_i)\); on odd edges let \(X_B\) be the weighted swap.
They vanish on unmatched endpoint basis vectors. Let \(W_B\) be the
unweighted swap on even edges, and \(W_A\) the unweighted swap on odd
edges, both with value \(+1\) on unmatched vectors. The matched blocks
satisfy \(ZX+XZ=0\) and \(Z^2+X^2=I\). At an unmatched endpoint
\(X=0\), \(Z=\pm I\), because the endpoint label is signed. Consequently
\[
 A_1=(I+Z_A-X_A)/2,\quad A_2=(I+Z_A+X_A)/2,\quad A_3=(I+W_A)/2
\]
and
\[
 B_1=(I+Z_B-X_B)/2,\quad B_2=(I+Z_B+X_B)/2,\quad B_3=(I+W_B)/2
\]
are orthogonal projectors. Expansion gives
\[
 \Bell=Z_AZ_B+\tfrac12 Z_A-\tfrac12 Z_B+
        \tfrac12 X_AW_B+\tfrac12 W_AX_B-I.
\]
For a real unit vector \(v\), its expectation on
\(\sum_i v_ie_i\otimes e_i\) is \(v^TJ_m(\gamma)v\):
the diagonal is \(d(\gamma_{i-1},\gamma_i)\) in either parity, and
each edge contributes \(w(\gamma_i)v_iv_{i+1}\). This proves the
finite realization used in Proposition~\ref{prop:qstar}. The source's
bilateral indexing is converted by \(c_j=\gamma_{j-1}\), with the
amplitude index unchanged. It gives exactly \eqref{eq:H}.

\subsection{A nonzero square-summable limit}\label{app:compactness}
For any finite path and real \(f\),
\begin{equation}\label{eq:TV}
 \langle f,J_m(\gamma)f\rangle
 \le \frac14\|f\|_2^2+\frac12\TV(\gamma)\|f\|_\infty^2.
\end{equation}
To see this, subtract \((x^2-1+y^2-1)/2\) from \(d(x,y)\):
the result is \(-(x-y)^2/2+(x-y)/2\le |x-y|/2\).
The remaining two-vertex form for an internal edge has matrix
\[
 \begin{pmatrix}-w^2/2&w/2\\w/2&-w^2/2\end{pmatrix},
\]
whose top eigenvalue is \(-w^2/2+w/2\le1/8\).
Summing counts each vertex at most twice; the signed endpoints
contribute zero. This proves \eqref{eq:TV}.

Take the monotone paths of External premise~\ref{prem:Coladangelo}(3)
and unit top eigenvectors \(f^{(n)}\ge0\). Absolute values cannot reduce
the Rayleigh quotient when off-diagonals are nonnegative, so these
nonnegative choices exist. Their values \(s_n\) tend to \(s>1/4\);
since their variation is two, \eqref{eq:TV} implies
\(\|f^{(n)}\|_\infty^2\ge s_n-1/4\).
Translate a maximum coordinate to zero. Extend the labels by \(+1\)
to the left and \(-1\) to the right, and amplitudes by zero. These are
bilateral eigenvectors because the separating endpoint edges have
zero off-diagonal entries.

A diagonal subsequence in the compact coordinate intervals gives
nonincreasing labels \(c_j\in[-1,1]\) and amplitudes \(f_j\ge0\).
Every finite sum of \(f_j^2\) is at most one, so \(f\in\ell^2(\ZZ)\).
The lower bound at zero makes \(f\ne0\). Continuity of \(b,d\) in each
three-term equation gives \(H(c)f=sf\). Set \(\lambda=f/\|f\|_2\).
At this point neither positivity at every site nor strict decrease
has been assumed.

\subsection{Boundary labels and support}
At a zero amplitude, the nonnegative eigenvalue equation forces
neighboring amplitudes across positive edges to vanish. Each connected
component of positive edges is therefore either identically zero or
strictly positive. Monotonicity of the labels leaves at most one
component containing positive edges. Other components are isolated
vertices with \(d\le0<s\), hence zero amplitudes. A nonzero component
exists and contains an edge.

Suppose this component has a left endpoint \(j_0\). Then
\(c_{j_0-1}=c_{j_0}=1\), \(\lambda_{j_0-1}=0\) and
\(\lambda_{j_0}>0\). Choose
\[
 0<\varepsilon<\min\{1,(1-3s)/2\},\qquad
 c'_{j_0}=1-\varepsilon,\qquad
 f=\lambda+\alpha\lambda_{j_0}e_{j_0-1},\qquad
 \alpha=\frac{\sqrt{2\varepsilon-\varepsilon^2}}{2s+\varepsilon}.
\]
A direct two-vertex computation gives
\[
 \frac{\langle f,H(c')f\rangle-s\|f\|_2^2}{\lambda_{j_0}^2}
 =\varepsilon\left[-(c_{j_0+1}+1/2)
                +\frac{2-\varepsilon}{4s+2\varepsilon}\right]>0.
\]
The last fraction exceeds \(3/2\), whereas \(c_{j_0+1}+1/2\le3/2\).
This contradicts the universal comparison in
Proposition~\ref{prop:qstar}. At a right endpoint \(j_1\), change
\(c_{j_1+1}=-1\) to \(-1+\varepsilon\) and add
\(\alpha\lambda_{j_1}e_{j_1+1}\). The bracket becomes
\(-(1/2-c_{j_1})+(2-\varepsilon)/(4s+2\varepsilon)>0\),
the same contradiction. Thus the component is all of \(\ZZ\):
\begin{equation}\label{eq:interiorpositive}
 -1<c_j<1,\qquad \lambda_j>0\quad(j\in\ZZ).
\end{equation}
This argument establishes the weaker positivity fact actually needed;
it does not use the false source sentence about all entries becoming
positive after a scalar identity shift.

\subsection{Finite label variations}
For fixed \(\lambda\), the Rayleigh form is
\[
 Q(c)=\sum_j d(c_j,c_{j+1})\lambda_j^2+
             \sum_j w(c_j)\lambda_{j-1}\lambda_j.
\]
It is absolutely convergent by bounded coefficients and
\(\sum_j\lambda_{j-1}\lambda_j\le1\). A finite change of interior
labels changes only finitely many smooth summands. The comparison
\(Q(c')\le s\), with equality at \(c\), makes the first derivative zero
and the Hessian nonpositive for every finite variation. In particular
\begin{equation}\label{eq:labelstationarity}
 (c_{j-1}-1/2)\lambda_{j-1}^2+
 (c_{j+1}+1/2)\lambda_j^2
 -\frac{c_j}{w(c_j)}\lambda_{j-1}\lambda_j=0
\end{equation}
and
\begin{equation}\label{eq:Hessian}
 \frac12D^2Q[\delta,\delta]
 =-\frac12\sum_j\frac{\lambda_{j-1}\lambda_j}{w(c_j)^3}\delta_j^2
          +\sum_j\lambda_j^2\delta_j\delta_{j+1}\le0.
\end{equation}
Together with the eigenvalue equation, \eqref{eq:labelstationarity}
gives \eqref{eq:stationarypair} and \eqref{eq:F}. All denominators
are positive by \eqref{eq:interiorpositive}.

\subsection{Equality geometry and exclusion of plateaux}
The function \(\mathfrak h\) in \eqref{eq:storage} is bounded, since setting
\(x=0\) gives \(\mathfrak h(y)\le s+1+y/2\); it is continuous in the interior
by concavity. Absolute convergence allows the deficit identity
\begin{align}
 0={}&s\|\lambda\|^2-Q(c)\label{eq:deficits}\\
 ={}&\sum_j[s-d(c_j,c_{j+1})-\mathfrak h(c_{j+1})-\mathfrak h(-c_j)]\lambda_j^2\notag\\
 &+\sum_j[\mathfrak h(c_j)\lambda_{j-1}^2+\mathfrak h(-c_j)\lambda_j^2
                  -2b(c_j)\lambda_{j-1}\lambda_j].\notag
\end{align}
Every summand is nonnegative. For an edge, this follows by completing
a square and using \(\mathfrak h(c_j)\mathfrak h(-c_j)\ge b(c_j)^2\). Since all amplitudes
are positive, equality forces
\begin{equation}\label{eq:saturation}
 \begin{gathered}
 \mathfrak h(c_{j+1})+d(c_j,c_{j+1})+\mathfrak h(-c_j)=s,\\
 \mathfrak h(c_j)\mathfrak h(-c_j)=b(c_j)^2,\qquad \mathfrak h(c_j)=b(c_j)u_j .
 \end{gathered}
\end{equation}

Here differentiability at saturated labels must be proved rather than
assumed. If \(t\in(-1,1)\) and \(\mathfrak h(t)\mathfrak h(-t)=b(t)^2\), both values of
\(\mathfrak h\) are positive. Concavity gives finite one-sided derivatives
\(\mathfrak h'_-\ge \mathfrak h'_+\) at both labels. For \(p(z)=\mathfrak h(z)\mathfrak h(-z)\),
\[
 p'_-(t)=\mathfrak h'_-(t)\mathfrak h(-t)-\mathfrak h(t)\mathfrak h'_+(-t),\qquad
 p'_+(t)=\mathfrak h'_+(t)\mathfrak h(-t)-\mathfrak h(t)\mathfrak h'_-(-t).
\]
Thus \(p'_--p'_+\) is the sum of two nonnegative derivative jumps
with positive weights. But \(p-b^2\ge0\) has a zero at \(t\), so
\(p'_-(t)\le-t/2\le p'_+(t)\). All inequalities are equalities, and
both jumps vanish. Consequently \(\mathfrak h\) is differentiable at \(t,-t\).
Varying the two arguments at a saturated pair \((x,y)\) gives
\begin{equation}\label{eq:predsucc}
 x=1/2-\mathfrak h'(y),\qquad y=\mathfrak h'(-x)-1/2.
\end{equation}
These formulas determine the predecessor and successor uniquely
among saturated equality pairs.

If \(c_j=c_{j+1}=t\), \eqref{eq:predsucc} propagates this equality
to each adjacent pair and hence to the entire sequence. For constant
labels the quadratic form has upper bound
\[
 t^2-1+w(t)=-w(t)^2+w(t)\le1/4
\]
on unit vectors, by Cauchy--Schwarz. This contradicts \(s>1/4\).
Therefore \(c_j>c_{j+1}\) everywhere.

\subsection{Interior tail limits}
Strict decrease gives distinct limits \(t_-,t_+\in[-1,1]\).
To exclude boundary limits before assuming ratio convergence, set
\(\mathsf A_j=w(c_j)u_j\) and \(\mathsf B_j=w(c_j)/u_j\).
Equation \eqref{eq:labelstationarity} and the Jacobi equation give
\begin{align}
 (c_{j+1}+1/2)\mathsf A_j+(c_{j-1}-1/2)\mathsf B_j&=c_j,
 &\mathsf A_j\mathsf B_j&=1-c_j^2,\label{eq:AB}\\
 \mathsf B_j+\mathsf A_{j+1}&=2(s-d(c_j,c_{j+1})).&&\notag
\end{align}
If a tail tends to \(+1\), the first equation bounds both positive
quantities on that tail. Their possible subsequential limits are
\((2/3,0)\) and \((0,2)\). If a tail tends to \(-1\), multiply the
first equation by \(-1\); the possible pairs are \((2,0)\) and
\((0,2/3)\). Simultaneous convergent subsequences at neighboring indices,
substituted in the second equation with \(d\to0\), would give
\(s\in\{0,1/3,1,4/3\}\). This contradicts the initial interval.
Hence both limits are interior, and the labels lie in a fixed compact
subinterval of \((-1,1)\).

\begin{lemma}[Positive Jacobi tail ratios]\label{lem:ratios}
Suppose a positive right-tail sequence satisfies
\[
 a_nf_{n-1}+d_nf_n+a_{n+1}f_{n+1}=sf_n,\qquad
 a_n\to a>0,\quad d_n\to d,\quad K=(s-d)/a>2.
\]
If it is square summable, then
\[
 \frac{f_n}{f_{n-1}}\longrightarrow
 \frac{K-\sqrt{K^2-4}}2<1.
\]
For a square-summable left tail the forward ratio tends to the
reciprocal root, which exceeds one.
\end{lemma}
\begin{proof}
Positivity gives, for \(v_n=f_n/f_{n-1}\),
\[
 v_n>\frac{a_n}{s-d_n},\qquad
 v_{n+1}<\frac{s-d_n}{a_{n+1}}.
\]
Thus late ratios stay in a compact positive interval. The recurrence is
\(v_{n+1}=F(v_n)+o(1)\), uniformly there, with \(F(v)=K-1/v\).
As \(F\) is increasing and continuous, the liminf and limsup are each
fixed points of \(F\). These are \(\rho_-<1<\rho_+\).
They cannot be different: for any \(m\in(\rho_-,\rho_+)\),
\(F(m)>m\), and sufficiently small late errors make crossing above
\(m\) irreversible. A limsup of \(\rho_+\) would force such a crossing,
contradicting liminf \(\rho_-\). The ratio therefore converges.
Square summability excludes \(\rho_+\), since ratios eventually bounded
below by a constant greater than one force growth. Reverse the sequence
for the left tail.
\end{proof}

For a carrier tail with label limit \(t\), the limiting coefficients
are \(a=w(t)/2\), \(d=t^2-1\). The condition \(K>2\) follows from
\[
 s>1/4\ge t^2-1+w(t).
\]
Lemma~\ref{lem:ratios} gives full positive finite ratio limits
\(u_->1\) on the left and \(u_+<1\) on the right.

\subsection{Tail elimination and exclusion of the extraneous branch}
Passing to the now justified limits in the two stationary equations,
write the label and ratio as \(t,r\), with \(w=\sqrt{1-t^2}>0\):
\begin{equation}\label{eq:tailequations}
 \frac w2(r+r^{-1})+t^2-1=s,\qquad
 (t+1/2)r+(t-1/2)r^{-1}=t/w.
\end{equation}
Putting \(h_0=s-t^2\) gives
\[
 r+r^{-1}=\frac{2(h_0+1)}w,\qquad
 r-r^{-1}=-\frac{2t(2h_0+1)}w.
\]
Their squared difference equals four, so
\begin{equation}\label{eq:tailpolynomial}
 0=(s-t^2)\{4t^4-(4s+5)t^2+s+2\}.
\end{equation}
The coefficient identity is checked exactly, without inserting a
preferred root branch.

In \eqref{eq:Hessian} take a finite variation
\(\delta_j=z_j/\sqrt{\lambda_{j-1}\lambda_j}\). Its finitely many
entries are finite even far out a small-amplitude tail, so a sufficiently
small perturbation parameter keeps the labels interior. Twice the
resulting form is
\[
 -\sum_j w(c_j)^{-3}z_j^2
       +2\sum_j\sqrt{u_j/u_{j+1}}\,z_jz_{j+1}\le0.
\]
Choose \(z=1\) on five consecutive sites and translate them down either
tail. The form tends to \(8-5/(1-t^2)^{3/2}\).
On \(s=t^2\), this is strictly positive, since
\[
 (1-s)^3>(37/50)^3>25/64=(5/8)^2.
\]
The last comparison is exact rational arithmetic. Sufficiently far
translated variations then contradict the nonpositive form. Hence
the branch \(s=t^2\) is excluded.

\subsection{Outer rather than inner or mixed tails}
The remaining polynomial \(P_s(z)=4z^2-(4s+5)z+s+2\) satisfies
\[
 P_s(1/4)=1,\qquad P_s(3/4)=1/2-2s<0,\qquad P_s(1)=1-3s>0.
\]
It has two roots with
\[
 1/4<\beta(s)^2<3/4<a(s)^2<1,\qquad
 \beta(s)=\sqrt{\frac{4s+5-\sqrt{16s^2+24s-7}}8}.
\]
Adding the sum and difference equations above and using
\[
 (2t+1)[(s+1-t^2)-t(2s+1-2t^2)]-(1+t)(2t-1)=P_s(t^2)
\]
gives \(r=r(t)\) as in \eqref{eq:r}. All denominators are nonzero
because the root magnitudes exceed \(1/2\). Moreover
\[
 r(t)^2-1=\frac{2t(4t^2-3)}{(1-t)(2t+1)^2},\qquad
 r(-t)=r(t)^{-1}.
\]
The left ratio \(>1\) permits labels \(a,-\beta\), and the right
ratio \(<1\) permits labels \(-a,\beta\).

Root ordering alone does not remove the same-sign pairs
\((a,\beta)\) and \((-\beta,-a)\). For this step use the scalar
certificate. Taking limits in \eqref{eq:saturation} gives
\[
 \mathfrak h(t)+\mathfrak h(-t)+t^2-1=s,\quad
 \mathfrak h(t)\mathfrak h(-t)=b(t)^2,\quad \mathfrak h(t)=b(t)r(t).
\]
The constant pair with both signs reversed also saturates
\eqref{eq:storage}. Differentiability and \eqref{eq:predsucc}
then imply, at either positive magnitude \(m=|t|\),
\[
 \mathfrak h'(m)=1/2-m,\qquad
 \mathfrak h(m)=f_0(m):=\frac{(1+m)(2m-1)}{2(2m+1)}
                 =\frac m2-\frac1{2(2m+1)}.
\]
If both magnitudes \(a>\beta>1/2\) occurred, the supporting line
at \(\beta\) would give
\[
 \mathfrak h(a)\le \mathfrak h(\beta)+(1/2-\beta)(a-\beta)<\mathfrak h(\beta).
\]
But
\[
 f_0(a)-f_0(\beta)
 =(a-\beta)\left(\frac12+\frac1{(2a+1)(2\beta+1)}\right)>0,
\]
a contradiction. Mixed magnitudes are impossible. Of the two
remaining equal-magnitude pairs, \((-\beta,\beta)\) is increasing,
so strict decrease leaves precisely \((a,-a)\).
The already established ratio limits are consequently
\(r(a)>1\) and \(r(-a)<1\). This proves all of
\eqref{eq:carrierconditions} at \(s=\St\), completing
Theorem~\ref{thm:R0}.

